\section{Results}
\label{sec:results}

\subsection{How to read the estimates}
\label{sec:results-read}
Three numbers appear throughout this section, and they answer different
questions.

\emph{The reduced form (RF)} is the raw causal quantity: how many
dollars the average credit-card tip changes at the cutoff, holding the
non-surcharge fare constant. It is the coefficient a rider would
recognize---``tips went up 29 cents at 4:00pm.''

\emph{ITT per statutory dollar} divides the reduced form by the
\emph{statutory} surcharge change (\$2.50, \$1.00, and so on). Because
the statute is a known constant, no first stage is estimated, no
exclusion restriction is invoked, and the standard error needs no
adjustment. It answers the policy question: for each dollar the rule
adds, how much do tips rise on average across \emph{all} standard-rate
trips---including any trip the meter failed to charge? It is the
design-based number, and the one we headline.

\emph{$\rho^{IV}$} instruments the measured surcharge with the cutoff
indicator and estimates both stages jointly, so its standard error
accounts for the first stage being estimated rather than known. With a
continuous treatment it is an average causal response
\parencite{angrist1995two}, and it is causal only under an exclusion
restriction: that the entire discontinuity in the measured surcharge
field is rush surcharge. Table~\ref{tab:composition} is why that
assumption cannot be verified---the field is contaminated on both
sides of every cutoff---so we read $\rho^{IV}$ as the behavioral rate
among trips whose recorded surcharge moved, corroborated externally by
the mechanical benchmark (Appendix~\ref{app:robustness}). Which scaling
matters depends on the question: pricing a fee \emph{before} enactment
uses ITT per statutory dollar; aggregating the transfer from fees
already collected uses $\rho^{IV}$, because collections are recorded
dollars. ITT sits below $\rho^{IV}$ exactly insofar as some trips'
recorded surcharge does not move at the cutoff.

One further habit of reading: every specification conditions on the
non-surcharge base (the metered fare and fixed fees). Section
\ref{sec:results-balance} shows why this is not optional.

\subsection{What else changes at 4:00pm: balance}
\label{sec:results-balance}
Table~\ref{tab:balance} reports discontinuities in trip characteristics
at the weekday 4:00pm cutoff. The honest news first: trips themselves
change. The metered fare jumps \$0.49, trips get 0.58 miles shorter and
0.69 minutes slower---the signature of rush-hour congestion. Continuity
of trip characteristics fails, which is why we describe the design as a
\emph{fare-conditioned} within-day discontinuity rather than a
conventional regression discontinuity: identification leans on
conditioning on the fare, and on the weekend contrast, not on an
assumption that nothing else moves at 4:00pm.

What does \emph{not} change is the composition of payers: the card share
moves 0.2 percentage points and the vendor mix 0.03 points. The people
tipping on either side of the cutoff are paying the same way, through
the same machines.

\begin{table}[!htbp]
  \centering
  \singlespacing
  \caption{\capttitle{Balance at the 4:00pm Cutoff}}
  \label{tab:balance}
  \small
  \begin{tabular}{@{}lccc@{}}
    \toprule
    Trip characteristic & Pre-cutoff mean & Jump at 4:00pm & (SE) \\
    \midrule
    \input output/tables/tab-balance.tex
    \bottomrule
  \end{tabular}
  \floatnotes{Weekday standard-rate card trips, \$2.50 regime, bandwidth
  $\pm 60$ minutes, local linear, clustered by date. The fare, distance,
  and duration jumps are why every estimate in this paper conditions on
  the non-surcharge base; the null card-share and vendor-share jumps are
  why payment composition is not a confound. \sym{***}, \sym{**},
  \sym{*}: significance at 1, 5, 10 percent.}
\end{table}

\subsection{The main estimates}
Figure~\ref{fig:4pm} shows the raw pattern the estimates formalize: tips
step up at 4:00pm in both regimes, by visibly more when the surcharge is
larger, then drift down through the evening as post-rush trips shorten
and cheapen---the fare dynamics the specification conditions away.
Figure~\ref{fig:wewd} shows the same clock on weekends, where no
surcharge exists: drift without a step.

\begin{figure}[!htbp]
  \centering
  \singlespacing
  \caption{\capttitle{Credit-Card Tips by Minute Around 4:00pm}}
  \label{fig:4pm}
  \includegraphics[width=0.86\textwidth]{fig-4pm.pdf}
  \floatnotes{Standard-rate card trips, weekdays excluding holidays,
  trip-weighted means by minute of pickup. Navy: \$1.00 regime. Orange:
  \$2.50 regime. The step at the dashed line is the treatment; the
  gradual decline afterward is rush-hour fare dynamics, absorbed by the
  fare control in estimation.}
\end{figure}

\begin{figure}[!htbp]
  \centering
  \singlespacing
  \caption{\capttitle{Weekdays Versus Weekends at the Same Clock}}
  \label{fig:wewd}
  \includegraphics[width=0.80\textwidth]{fig-wewd.pdf}
  \floatnotes{Mean credit-card tips by minute around 4:00pm, \$2.50
  regime, standard-rate card trips. Weekdays (navy) receive the rush
  surcharge from 4:00pm; weekends (teal) never do.}
\end{figure}

Table~\ref{tab:main3} is the paper's central table: the same
specification at all six within-day surcharge changes, with the two
weekend placebos beneath. Reading across the top row: at the weekday
4:00pm cutoff in the \$2.50 regime, tips rise \$0.292 (SE \$0.005); per
statutory dollar that is $0.117$; instrumenting the measured surcharge,
$\rho^{IV} = 0.150$ (SE $0.002$). Reading down the $\rho^{IV}$ column is
the paper's result: $0.150$, $0.141$, $0.134$, $0.134$, $0.122$,
$0.128$. Two doses, two clocks, both directions, one parameter near
$0.14$. The placebo rows show what the same machinery finds when no
surcharge changes: \$0.024 and $-$\$0.002.

\begin{table}[!htbp]
  \centering
  \singlespacing
  \caption{\capttitle{Pass-Through at Six Surcharge Changes}}
  \label{tab:main3}
  \footnotesize
  \setlength{\tabcolsep}{4.5pt}
  \begin{tabular}{@{}lccccc@{}}
    \toprule
    & \multicolumn{2}{c}{Reduced form (USD)} & ITT per & \multicolumn{2}{c}{$\rho^{IV}$} \\
    \cmidrule(lr){2-3} \cmidrule(lr){5-6}
    Surcharge change & Tip step & (SE) & statutory \$ & Estimate & (SE) \\
    \midrule
    \input output/tables/tab-main3.tex
    \bottomrule
  \end{tabular}
  \floatnotes{Standard-rate card trips, September 2021--September 2024,
  bandwidth $\pm 60$ minutes, local linear in pickup minute, conditioning
  on the non-surcharge base, weighted by card trips, clustered by date.
  ITT divides the reduced form by the statutory dose (a known constant,
  so its SE is the RF's SE rescaled). $\rho^{IV}$ instruments the
  measured card-trip surcharge with the cutoff indicator; its SE reflects
  joint estimation of both stages. Weekend placebo rows apply the
  identical specification where no surcharge changes. \sym{***},
  \sym{**}, \sym{*}: significance at 1, 5, 10 percent.}
\end{table}

Figure~\ref{fig:coefplot} is the replication argument in one frame:
every locked change, expressed per statutory dollar, on one axis. The
six primary changes span $0.090$--$0.117$ with a precision-weighted
mean of $0.108$ (SE $0.001$). The 6:00am weekend changes sit at
$0.126$ and $0.155$ with wide intervals. The JFK flat-fare surcharge
(squares; Section~\ref{sec:results-controls}) contributes the \$4.50
and \$5.00 doses at $0.133$ and $0.149$. The 2025 congestion toll,
estimated with no clock at all, lands at $0.107$. The weekend placebos
sit at zero, and the two 6:00am weekday failures appear in grey rather
than disappearing (Section~\ref{sec:results-6am}).

\begin{figure}[!htbp]
  \centering
  \singlespacing
  \caption{\capttitle{One Parameter: Pass-Through per Statutory Dollar
  at Every Locked Change}}
  \label{fig:coefplot}
  \includegraphics[width=0.94\textwidth]{fig-coefplot.pdf}
  \floatnotes{Intention-to-treat tip response per statutory surcharge
  dollar, with 95 percent confidence intervals; specification of
  Table~\ref{tab:main3} throughout. Navy circles: within-day surcharge
  changes on the standard-rate sample. Navy squares: the JFK flat-fare
  surcharge (Section~\ref{sec:results-controls}), whose higher ITT
  reflects near-complete treatment delivery (94 versus 78 percent of
  the statutory dose), not different behavior. Grey circles: the 6:00am
  weekday changes, where the dawn compositional shift breaks the fare
  conditioning (Section~\ref{sec:results-6am}). Teal diamond: the
  January 2025 congestion toll, estimated from the spatial
  origin--destination design with no clock involved, divided by the
  \$0.75 statutory toll. Hollow circles: weekend placebos at 4:00pm,
  scaled by the matched regime's dose. The dashed line marks the
  precision-weighted mean of the six primary standard-rate changes,
  $0.108$ (SE $0.001$). Minimum detectable effects (2.8 $\times$ SE)
  for the six primary changes range from $0.005$ to $0.039$ per
  statutory dollar.}
\end{figure}

Estimated as an explicit difference-in-discontinuities---the weekday
step minus the weekend step at the same clock
\parencite{grembi2016difference}, the specification of
equation~\eqref{eq:didisc}---the effect is \$0.271 (SE \$0.008) in the
\$2.50 regime and \$0.105 (\$0.008) in the \$1.00 regime: $0.108$ and
$0.105$ per statutory dollar. Dropping the fare control entirely moves
these to $0.120$ and $0.092$. The conclusion does not lean on the
functional form of the fare adjustment; the weekend difference carries
it.

The fuzzy version of the same design reports all three quantities of
the standard presentation. The first stage---the weekday--weekend
difference in the \emph{surcharge} discontinuity---is \$1.934 (SE
\$0.011) in the \$2.50 regime; the reduced form is the \$0.271 above;
and the just-identified 2SLS coefficient, instrumenting the measured
surcharge with the cutoff-by-weekday interaction, is
$\hat{\tau} = 0.140$ (SE $0.004$). In the \$1.00 regime the same three
numbers are \$0.713 (\$0.007), \$0.105 (\$0.008), and $0.147$
($0.011$): the fuzzy estimate is dose-stable, and it sits inside the
$\rho^{IV}$ family, as it must, since it differs from the single-cutoff
IV only in differencing the weekend counterfactual through both
stages.

\subsection{A working-day control, and a \$5 dose on a fare that cannot move}
\label{sec:results-controls}
Two further designs close the remaining doors on the 4:00pm result, and
both were run after the main estimates as out-of-sample checks.

\emph{Holidays.} Legal holidays are calendar weekdays on which the rush
surcharge does not apply---a working-day counterfactual that weekends
cannot supply. The federal holiday list is not the TLC list, so we
classify each holiday by its own first stage: the traditional holidays
are cleanly exempt, one observed date is partially charged and dropped,
and Juneteenth is \emph{fully charged} in all three years (first stages
of \$0.90, \$2.29, and \$2.28 against statutory doses of \$1.00 and
\$2.50)---the meters do not observe the newest federal holiday.
Differencing the weekday step against the exempt holidays gives $0.096$
per statutory dollar in the \$2.50 regime and $0.095$ in the \$1.00
regime (wild-cluster bootstrap $p < 10^{-4}$ and $p = 0.035$; holiday
clusters are few, so the bootstrap is the relevant inference)---the
same answer the weekend difference gives. And Juneteenth supplies the
converse check: a day with holiday \emph{composition} but the full
\emph{dose} shows a tip step of \$0.241 (SE \$0.067), $0.096$ per
statutory dollar. Composition off, dose on: the step follows the dose.

\emph{JFK flat fares.} The analysis plan quarantines special rate codes
from the main sample; analysed separately, the JFK flat fare is the
cleanest design in the paper. Its rush-hour surcharge is \$4.50 in the
early regime and \$5.00 after December 2022---twice the largest dose in
the main design---and it sits on a \emph{flat} fare, so the metered fare
cannot jump at 4:00pm by construction: the confound that motivates the
fare conditioning everywhere else does not exist here. The estimates
behave accordingly. The tip step is \$0.745 (SE \$0.052) at the \$5.00
dose and \$0.596 (\$0.052) at \$4.50---$0.149$ and $0.133$ per
statutory dollar---and dropping every control changes the answer by
three hundredths of a cent ($0.1486$ against $0.1489$). Weekend
placebos are null on both sides (tip $-\$0.03$ and $-\$0.05$, first
stages $-0.01$ and $-0.04$). Per \emph{recorded} surcharge dollar the
JFK estimates are $0.159$ and $0.150$, on top of the main design's
$\rho^{IV}$; the higher ITT relative to the pooled $0.108$ is treatment
delivery, not behavior---JFK meters deliver 94 percent of the statutory
dose against 78 percent on standard rates. The mechanical benchmark
tells the same story at five times the stakes: the pre-4:00pm JFK tip
rate of $17.5$ percent predicts a step of \$0.875 under complete
pass-through, and the estimate is 85 percent of it, the same fraction
as the main design. Riders tipping on an \$85 airport fare pass a \$5
surcharge into the tip at the same rate riders tipping on a \$25
Manhattan fare pass \$2.50.

Interpreting the magnitude: the average tip in this sample is 17--18
percent of the base. If riders treated a surcharge dollar exactly like a
fare dollar, complete default-mediated pass-through would put $\rho$
near $0.17$; if riders fully distinguished the surcharge---tipping on
the service, not the stack---$\rho$ would be $0$. The estimate of
$0.14$--$0.15$ says riders pass surcharges into tips at about
four-fifths of the rate they tip the fare itself: not total blindness to
the base, but close.

\subsection{Robustness}
Table~\ref{tab:bw3} shows $\rho^{IV}$ across bandwidths of 30 to 90
minutes ($0.147$--$0.156$) and donuts excluding 3, 5, or 10 minutes
around the cutoff ($0.148$--$0.152$). Because the surcharge is assessed
on the pickup time, the donut is not needed for assignment reasons; we
report it because cells adjacent to the cutoff mix minutes unevenly
within trips, and it makes no difference. Placebo cutoffs at 3:15pm and
5:00pm give tip-side effects of $-$\$0.0003 (SE \$0.007) and $-$\$0.008
(SE \$0.007)---both sides of the placebo, first stage and tip, are now
reported and both are null. A permutation over 60 candidate cutoffs,
each estimated at the same $\pm 30$ minute bandwidth as the true cutoff
and drawn only from windows that never cross 4:00pm, finds none with an
effect as large as the true cutoff's, and a wild cluster bootstrap on
the corrected reduced form rejects zero at $p < 10^{-6}$. Minimum
detectable effects---2.8 times the standard error, the analysis plan's
locked criterion---range from $0.005$ to $0.039$ per statutory dollar
across the six primary changes, so even the least-powered design could
detect an effect one-third the size of the pooled estimate.

\begin{table}[!htbp]
  \centering
  \singlespacing
  \caption{\capttitle{Bandwidth and Donut Robustness, $\rho^{IV}$}}
  \label{tab:bw3}
  \small
  \begin{tabular}{@{}lcc@{}}
    \toprule
    Specification & $\rho^{IV}$ & (SE) \\
    \midrule
    \input output/tables/tab-bw3.tex
    \bottomrule
  \end{tabular}
  \floatnotes{Weekday 4:00pm cutoff, \$2.50 regime, corrected
  specification, clustered by date. \sym{***}, \sym{**}, \sym{*}:
  significance at 1, 5, 10 percent.}
\end{table}

\subsection{The third clock: 6:00am, and where the design breaks}
\label{sec:results-6am}
The analysis plan locks a third clock. At 6:00am the overnight surcharge
\emph{ends}, so the statutory step is downward---\$0.50 before December
2022, \$1.00 after---and because the overnight surcharge applies on
every day of the week, there is no weekend placebo at this cutoff. That
absence turns out to be informative.

On weekends the estimates behave: $\rho^{IV} = 0.175$ (SE $0.041$) at
the \$1.00 dose and $0.165$ ($0.088$) at \$0.50, a little above the
main cluster but within its neighborhood. On weekdays they do not: the
\$1.00 dose gives $0.016$ ($0.018$) and the \$0.50 dose gives $-0.143$
($0.037$)---null and wrong-signed respectively.

We report this because the plan locked it, and because the pattern
diagnoses itself. The same clock and the same statutory dose produce
sensible estimates on weekends and broken ones on weekdays. What differs
between them is what happens at dawn on a working day: the pickup mix
shifts sharply toward airport-bound and commuter trips, which are longer
and dearer in a way that conditioning on the mean non-surcharge fare
does not absorb. On weekends that shift is absent and the design
recovers the same parameter as the 4pm and 8pm cutoffs. The honest
reading is that this design identifies pass-through wherever rider
composition evolves smoothly through the cutoff, and that the 6:00am
weekday window is the one place in our data where it plainly does not.
We therefore exclude the two weekday 6am estimates from the headline
range and report them here rather than in Table~\ref{tab:main3}.

\subsection{Do riders undo it?}
At the 4:00pm cutoff the share of tips exactly at a default button falls
0.85 percentage points, zero-tipping rises 0.15 points, and card use
\emph{rises} 0.19 points. Against a mechanical benchmark---the surcharge
step times the prevailing 17.4 percent tip rate---the estimated
pass-through is about nine-tenths of complete. Riders offset almost none
of the base change.

\subsection{The button is not the mechanism: vendor evidence}
\label{sec:results-vendor}
The two payment vendors offer a built-in test of \emph{how} the
pass-through operates, and it refutes our own initial hypothesis. If
pass-through worked purely through riders pressing default buttons, it
should scale with button use: Vendor 1's riders take an exact default on
19 percent of tips, Vendor 2's on 65 percent, so $\rho$ should differ
by a factor of three. It does not. Estimated separately on the same 4pm
design, $\rho^{IV} = 0.147$ (SE $0.005$) at Vendor 1 and $0.152$
($0.003$) at Vendor 2---statistically indistinguishable. Riders who
type custom amounts evidently compute their percentages on the same
fee-inclusive total the screen displays. The operative mechanism is the
\emph{base the interface presents}, not the button that presents it,
which broadens the result: redesigning buttons would not undo the
fiscal channel; redesigning the base would. Figure~\ref{fig:vendor}
shows the contrast in one frame: adherence bars that differ by a factor
of 3.5, pass-through dots that coincide.

\begin{figure}[!htbp]
  \centering
  \singlespacing
  \caption{\capttitle{Same Pass-Through, Very Different Button Use}}
  \label{fig:vendor}
  \includegraphics[width=0.76\textwidth]{fig-vendor.pdf}
  \floatnotes{Bars (right axis): share of card tips exactly at a
  default button, pre-cutoff, by payment vendor. Dots (left axis, 95
  percent CIs): $\rho^{IV}$ from the weekday 4:00pm design estimated
  separately by vendor, \$2.50 regime. If pass-through operated only
  through button presses, the dots would differ as much as the bars.}
\end{figure}

\subsection{Who passes it through}
\label{sec:results-decomp}
The mechanism has a sharper test than vendor comparisons: classify every
tip by \emph{how it was chosen}, and ask whose tips move at the cutoff.
We split card tips into four mutually exclusive types---exactly a
20/25/30 percent button, a whole-dollar amount, other custom amounts,
and zero---and estimate the weekday--weekend difference-in-
discontinuities separately for each type's mean tip and its share of
trips (Table~\ref{tab:decomp}).

If pass-through were mechanical rounding---riders nudging totals to the
nearest dollar---whole-dollar tips would carry it. They do not. In the
\$2.50 regime, button tips rise \$0.421 (SE \$0.011) and custom
non-round tips \$0.247 (\$0.027), while whole-dollar tips move \$0.043
(\$0.025). The dose comparison is more damning still: when the statutory
dose was \$1.00, button tips rose \$0.182---the response scales
one-for-one with the legislated dose---while the whole-dollar response
is \$0.04 in \emph{both} regimes, insensitive to a 150 percent dose
change. Riders who think in percentages of the displayed total pass the
surcharge through; riders who think in round dollar amounts are inert.
And nobody switches type at the cutoff: every share moves by less than
half a percentage point. The aggregate offset from complete mechanical
pass-through is not partial attention spread thinly---it is the
whole-dollar fifth of riders sitting still.

\begin{table}[!htbp]
  \centering
  \singlespacing
  \caption{\capttitle{Who Passes It Through: DiDisc by Tipping Type}}
  \label{tab:decomp}
  \small
  \begin{tabular}{@{}lccc@{}}
    \toprule
    & DiDisc & (SE) & Per statutory \$ \\
    \midrule
    \input output/tables/tab-decomp.tex
    \bottomrule
  \end{tabular}
  \floatnotes{Weekday--weekend difference-in-discontinuities at 4:00pm,
  \$2.50 regime, fare-conditioned, clustered by date; tip rows weight by
  the category's own trips, share rows by card trips. Types are mutually
  exclusive: button = tip equals 20/25/30 percent of the pre-tip total
  to within half a cent; whole-dollar = round-dollar tips not at a
  button. In the \$1.00 regime the button response is \$0.182 (0.009)
  and the whole-dollar response \$0.043 (0.027): button pass-through
  scales with the dose, whole-dollar tipping is inert at both doses.}
\end{table}

\subsection{Demand: riders do not dodge, and barely notice}
\label{sec:results-demand}
Three further margins complete the behavioral picture, and all three are
near-zero in the same direction.

\emph{No retiming.} If riders timed hails to dodge the surcharge, pickup
density would pile up before 4:00pm and hollow out after. It does not:
the share of window pickups per minute is essentially identical across
the cutoff (ten-minute ratio $0.997$ at the \$1.00 dose, $0.987$ at
\$2.50; Figure~\ref{fig:density4pm}). A 150 percent increase in the fee
produced a one percent change in timing. This doubles as the design's
sorting diagnostic: the running variable is not manipulated.

\emph{No demand response at the toll.} Log eligible trips at the January
2025 congestion boundary move $-1.0$ percent (SE $0.9$)---statistically
and economically indistinguishable from zero for a \$0.75 fee on a
\$27 ride. Combined with the tip pass-through, the driver-side ledger of
the toll is striking: drivers gain roughly \$2.2 million a year in
induced tips on eligible card trips and lose no detectable fare revenue.
At the point estimate, the congestion toll made drivers net
\emph{better} off; the confidence interval cannot exclude a net loss if
the true demand decline sits at its lower bound.

\emph{No reaction to a 150 percent surcharge increase.} On December 19,
2022 the rush surcharge jumped from \$1.00 to \$2.50 overnight. If
riders monitored their totals, percentage tipping should have dropped by
about $1.1$ points in the rush window (the fall required to keep tip
dollars fixed). Netting seasonal drift with the surcharge-free pre-rush
window, the tip rate fell $0.22$ points (SE $0.07$)---riders undid about
a fifth of the increase and passed the rest through
(Figure~\ref{fig:dec22}). That one-fifth offset independently implies
pass-through of roughly four-fifths of the prevailing tip rate, which is
precisely where the clock designs put $\rho$. And the offset is not the
start of slow learning: Figure~\ref{fig:lp-dec22} traces the response
horizon by horizon, local-projection style, for six months after the
repricing. The path settles immediately at about $-0.2$ points and
never drifts toward the $-1.1$-point full-attention benchmark. Riders
who were going to notice noticed at once; the rest never did.

\begin{figure}[!htbp]
  \centering
  \singlespacing
  \caption{\capttitle{No Delayed Learning: The Adaptation Path After December 2022}}
  \label{fig:lp-dec22}
  \includegraphics[width=0.82\textwidth]{fig-lp-dec22.pdf}
  \floatnotes{Weekly horizon-by-horizon responses of the rush-window
  tip rate, netted against the surcharge-free pre-rush window and
  measured relative to the 26-week pre-period mean, for 26 weeks after
  the \$1.00-to-\$2.50 repricing. The shaded band is a 95 percent
  interval from daily variation within each week. The maroon line marks
  the response that would keep tip dollars unchanged; zero marks full
  pass-through.}
\end{figure}

\begin{figure}[!htbp]
  \centering
  \singlespacing
  \caption{\capttitle{Pickup Density Through the 4:00pm Cutoff}}
  \label{fig:density4pm}
  \includegraphics[width=0.78\textwidth]{fig-density4pm.pdf}
  \floatnotes{Each date's share of window pickups in each minute,
  averaged by minute, weekdays. Smooth density through the cutoff means
  riders do not retime hails around the surcharge and the running
  variable is unmanipulated.}
\end{figure}

\begin{figure}[!htbp]
  \centering
  \singlespacing
  \caption{\capttitle{The December 2022 Surcharge Increase: Did Anyone Notice?}}
  \label{fig:dec22}
  \includegraphics[width=0.80\textwidth]{fig-dec22.pdf}
  \floatnotes{Weekly mean credit-card tips as a share of the pre-tip
  total, weekdays, in the rush window (navy; surcharge rose \$1.00 to
  \$2.50 at the dashed line) and the surcharge-free pre-rush window
  (teal). Keeping tip dollars fixed would require the navy series to
  fall about 1.1 points relative to teal; it falls 0.22.}
\end{figure}

\subsection{Transportability: congestion fees}
\label{sec:results-fee}
The clock designs identify $\rho$ from surcharges that switch on a
schedule; congestion policy switches them on at boundaries in space.
The re-estimated design assigns treatment by the statutory rule---a trip
is eligible if its pickup \emph{or} dropoff lies in the fee zone---on
standard-rate card trips, and reports joint pre-trend tests. At the
January 2025 toll the origin--destination design gives a pooled tip
effect of \$0.080 (SE \$0.005) against a first stage of \$0.602 of
fee per card trip: $\rho = 0.133$ (SE $0.009$), within a cent of the
clock designs. Correcting the eligibility matters exactly as
contamination predicts: the old pickup-only assignment gave \$0.058,
because treated dropoffs sat in its control group, and its apparent
post-event fade disappears under the correct rule---post coefficients
are stable at \$0.041--\$0.052 across all six months. Two refinements
discipline the pre-period. First, seasonality: December's holiday
tipping makes any single reference month misleading, so we add
treated-by-calendar-month fixed effects, identified from a second
pre-year (2023), and measure post effects against all pre months. The
effects strengthen and stabilize---\$0.074 to \$0.089 in every one of
the six post months ($t > 8$; Figure~\ref{fig:es-fee25})---but a joint
Wald test on the 2024 pre-year against 2023 still rejects
($p = 0.0007$): treated zones genuinely drifted during 2024's recovery,
beyond seasonality. Second,
and more decisively, the eligibility rule itself supplies a
\emph{continuous dose}: the post-period share of each pickup--dropoff
pair's trips actually charged the fee, which nets out the exempt
highway corridors no geometric rule can see. Conditioning on the
non-surcharge base---the same correction the clock designs
required---the dose design yields a monotone staircase: pairs charged
up to half the time show \$0.007 (SE \$0.011), pairs charged 50--90
percent show \$0.061 ($0.015$), and pairs charged above 90 percent
show \$0.085 ($0.013$), for an implied $\rho = 0.121$. Effects that
rise step-for-step with measured treatment intensity
(Figure~\ref{fig:dose}) are difficult to attribute to zone-level
drift, which would have to reproduce the dose gradient by coincidence.

\begin{figure}[!htbp]
  \centering
  \singlespacing
  \caption{\capttitle{The Dose Staircase: Tip Effects Rise with Charging Intensity}}
  \label{fig:dose}
  \includegraphics[width=0.76\textwidth]{fig-dose.pdf}
  \floatnotes{Post-toll tip effects by the share of each
  pickup--dropoff pair's trips actually charged the 2025 congestion
  fee (measured from post-period fee incidence), relative to pairs
  charged up to 10 percent of the time. Fare-conditioned, OD-pair and
  month fixed effects, clustered by pickup zone; bars are 95 percent
  confidence intervals.}
\end{figure} Stacking both fee adoptions into one design---each cohort with its own
unit and time effects, tips regressed on the observed fee instrumented
by eligibility-times-post, fare-conditioned throughout---yields the
\emph{average rollout effect} of putting a fee on the meter:
\[
  \bar{\rho}_{\text{rollout}} = 0.163 \ (\text{SE } 0.021).
\]
Within the same specification the cohorts are $0.171$ ($0.033$) for
2019 and $0.138$ ($0.007$) for 2025, statistically indistinguishable
($p = 0.36$). The uniform fare conditioning also resolves this
section's one outlier: the earlier 2019 estimate of $0.218$ carried
the same omitted-fare bias the clock designs taught us to remove, and
falls into line once treated identically. We still keep the spatial
evidence in a corroborating role---the binary event studies do not
pass a strict pre-trend test and we say so---while noting that the
dose gradient and the pooled rollout estimate land within two cents of
the clock designs.

\begin{table}[!htbp]
  \centering
  \singlespacing
  \caption{\capttitle{Pass-Through at Congestion Fees: The Initial
  Pickup-Only Design}}
  \label{tab:fee}
  \small
  \begin{tabular}{@{}lcccc@{}}
    \toprule
    Policy & First stage & Tip effect & (SE) & $\hat{\rho}$ \\
    \midrule
    \input output/tables/tab-fee.tex
    \bottomrule
  \end{tabular}
  \floatnotes{Zone-by-fare-bin difference-in-differences, zone-fare and
  month fixed effects, clustered by pickup zone. Windows:
  2018m7--2019m12 and 2024m1--2025m6. \emph{This is the initial
  pickup-only specification and is superseded}: it assigns treatment by
  pickup zone alone, so trips whose dropoff was in the fee zone sit in
  its control group. It is retained to show what the contamination does.
  The preferred estimates---origin--destination eligibility, fare
  conditioning, and the continuous dose---are reported in the text and
  give $\rho$ of $0.133$, $0.121$, and $0.163$ rather than the $0.146$
  and $0.254$ shown here.}
\end{table}

\begin{figure}[!htbp]
  \centering
  \singlespacing
  \caption{\capttitle{Event Study: The January 2025 Congestion Toll}}
  \label{fig:es-fee25}
  \includegraphics[width=0.78\textwidth]{fig-es-fee25.pdf}
  \floatnotes{Corrected design: origin--destination eligibility with
  treated-by-calendar-month fixed effects identified from calendar 2023;
  the twelve 2024 pre-months are shown as deviations from that seasonal
  baseline, and post coefficients are stable at \$0.074--\$0.089 with
  $t > 8$ throughout. The pre-period looks flat, and no single 2024
  coefficient is large---but a joint Wald test on all twelve rejects
  ($p = 0.0007$): treated routes drifted slightly during 2024's
  recovery. That is why the spatial evidence is read as corroboration,
  anchored by the dose gradient of Figure~\ref{fig:dose}, rather than
  as the headline. Bars: 95 percent confidence intervals.}
\end{figure}

\subsection{The size and geography of the transfer}
At $\rho = 0.15$, the \$127.2 million of 2019 congestion surcharges on
card trips implies roughly \$18.7 million of additional gratuities, and
the 2025 toll's \$17.4 million annualized implies about \$2.2 million a
year. Figure~\ref{fig:map-causal} is the spatial design's summary exhibit:
the per-zone, fare-conditioned change in tips---measured relative to
never-eligible zones---estimated at a placebo date (July 2024) and at
the real toll (January 2025), with zones muted unless significant at
five percent. At the placebo date, 6 percent of non-CBD zones and 10
percent of CBD zones show anything, scattered at random. At the real
date, \textbf{82 percent of congestion-district zones turn
significantly positive} (mean $+\$0.084$) against 18 percent
elsewhere: the tolled district materializes from the estimates alone,
with no boundary drawn. Figure~\ref{fig:map} then maps where the
resulting transfer is generated:
congestion-district pickup zones average \$46{,}000 of induced tips a
year each, against \$900 outside, so the policy boundary is visible in
the map without being drawn.

\begin{figure}[!htbp]
  \centering
  \singlespacing
  \caption{\capttitle{The Toll Appears in the Map: Placebo Versus Real Date}}
  \label{fig:map-causal}
  \includegraphics[width=0.98\textwidth]{fig-map-causal.pdf}
  \floatnotes{Per-pickup-zone change in mean credit-card tips,
  conditioned on fare bins and dropoff eligibility and measured relative
  to never-eligible zones, at a placebo split (July 2024, left) and the
  actual toll (January 2025, right). Blues: significant increases
  (darker larger); orange: significant decreases; grey: no significant
  change or insufficient trips. The congestion district emerges only at
  the real date.}
\end{figure}

\begin{figure}[!htbp]
  \centering
  \singlespacing
  \caption{\capttitle{Where the Transfer Is Generated}}
  \label{fig:map}
  \includegraphics[width=0.80\textwidth]{fig-map.pdf}
  \floatnotes{Annualized additional tip dollars induced by the 2025
  congestion toll, by pickup taxi zone: fee dollars charged on card
  trips times the estimated pass-through. Darker = larger; grey = none.}
\end{figure}
